\section{The PANOPTES Reference Architecture}
\label{sec:panoptes_architecture}

To address the diagnostic dilemmas identified in Section~\ref{sec:background} while strictly observing Control Principles \textbf{CP1} and \textbf{CP2}, we formalize the \textbf{PANOPTES} reference architecture. As illustrated in Figure~\ref{fig:panoptes_architecture} and categorized in Table~\ref{tab:panoptes_layers}, PANOPTES decomposes the cloud-native observability lifecycle into four decoupled, synergistic layers.

\begin{figure*}[!htbp]
\centering
\begin{tikzpicture}[
    font=\sffamily\scriptsize,
    basebox/.style={draw, rounded corners=3pt, align=center, inner sep=3.5pt},
    layerbox/.style={draw=black!60, rounded corners=4pt, dashed, inner sep=6pt, fill=black!2},
    titlelabel/.style={font=\sffamily\bfseries\footnotesize, text=black!85, anchor=west},
    sutnode/.style={basebox, fill=blue!10, draw=blue!65!black, minimum width=3.4cm, minimum height=0.62cm},
    probenode/.style={basebox, fill=cyan!12, draw=cyan!75!black, minimum width=5.0cm, minimum height=0.72cm},
    storenode/.style={basebox, fill=teal!12, draw=teal!75!black, minimum width=5.0cm, minimum height=0.75cm},
    correnode/.style={basebox, fill=purple!12, draw=purple!80!black, minimum height=0.85cm},
    actnode/.style={basebox, fill=orange!12, draw=orange!80!black, minimum width=11.2cm, minimum height=0.70cm},
    arrow/.style={-{Stealth[scale=0.85]}, thick, draw=black!75},
    biarrow/.style={{Stealth[scale=0.8]}-{Stealth[scale=0.8]}, thick, draw=black!70}
]

    % =========================================================================
    % LAYER 1: SUT & TELEMETRY INGESTION (PROBES)
    % =========================================================================
    % Microservices SUT Row
    \node[sutnode] (sut_frontend) at (-5.6, 6.2) {\textbf{frontend} / \textbf{cart}\\(HTTP/gRPC Gateway \& State)};
    \node[sutnode] (sut_shipping) at (0.0, 6.2) {\textbf{shipping} / \textbf{checkout}\\(Synchronous RPC Workloads)};
    \node[sutnode] (sut_payment)  at (5.6, 6.2) {\textbf{payment} / \textbf{currency}\\(Kafka Async \& Background)};

    % Probe Row (Aligned vertically with the 3 pillars)
    \node[probenode] (probe_traces) at (-5.6, 4.9) {
        \textbf{OTel SDK \& Auto-Instrumentation}\\
        \textit{Modality: Distributed Traces}\\
        W3C TraceContext In-Band Propagation
    };

    \node[probenode] (probe_metrics) at (0.0, 4.9) {
        \textbf{eBPF Probes \& cgroup v2 Scrapers}\\
        \textit{Modality: Infrastructure \& Socket Metrics}\\
        \texttt{kprobe/tcp\_v4\_connect} $\;|\;$ CPU Throttling
    };

    \node[probenode] (probe_logs) at (5.6, 4.9) {
        \textbf{Promtail / Fluentbit DaemonSet}\\
        \textit{Modality: Structured Container Logs}\\
        \texttt{/var/log/pods/*.log} Regex Streaming
    };

    \begin{scope}[on background layer]
        \node[layerbox, fit=(sut_frontend)(sut_payment)(probe_traces)(probe_logs)] (layer1) {};
    \end{scope}
    \node[titlelabel] at (layer1.north west) [yshift=0.25cm] {\textbf{Layer 1: Telemetry Ingestion \& Non-Invasive Kernel/User Probes}};

    % =========================================================================
    % LAYER 2: MULTIMODAL FEDERATED STORAGE BACKENDS
    % =========================================================================
    \node[storenode] (store_tempo) at (-5.6, 3.1) {
        \textbf{Grafana Tempo Object Store}\\
        \textit{Modality: Distributed Traces}\\
        OTLP Ingestion \& Inverted Trace Indexing
    };

    \node[storenode] (store_prom) at (0.0, 3.1) {
        \textbf{Prometheus TSDB}\\
        \textit{Modality: Time-Series Metrics}\\
        High-Frequency Scraping ($15\,\text{s}$) \& PromQL
    };

    \node[storenode] (store_loki) at (5.6, 3.1) {
        \textbf{Grafana Loki Chunk Store}\\
        \textit{Modality: Structured Event Logs}\\
        Label-Indexed Stream Storage \& LogQL
    };

    \begin{scope}[on background layer]
        \node[layerbox, fit=(store_tempo)(store_loki)] (layer2) {};
    \end{scope}
    \node[titlelabel] at (layer2.north west) [yshift=0.25cm] {\textbf{Layer 2: Multimodal Federated Telemetry Backends}};

    % =========================================================================
    % LAYER 3: CAUSAL CORRELATION & DIAGNOSTIC REASONING ENGINE
    % =========================================================================
    \node[correnode, minimum width=5.6cm] (algo_selftime) at (-5.3, 1.2) {
        \textbf{Algorithm 1: Span Self-Time Decomposition}\\
        {\footnotesize $T_{\text{self}} = T_{\text{span}} - \sum_{k} T_{\text{child}_k}$}\\
        Eliminates Cascading Span Bias on Ingress RPCs
    };

    \node[correnode, minimum width=9.8cm] (algo_fusion) at (2.8, 1.2) {
        \textbf{Algorithm 2: Cross-Modal Evidence Fusion}\\
        {\footnotesize $\mathcal{V}^* = \arg\max_{s} \Big(\omega_m \mathcal{S}_{\text{metric}} + \omega_l \mathcal{S}_{\text{log}} + \omega_t \mathcal{S}_{\text{trace}}\Big)$}\\
        Correlates Exemplar Trace Context with Metric/Log Signatures
    };

    \begin{scope}[on background layer]
        \node[layerbox, fit=(algo_selftime)(algo_fusion)] (layer3) {};
    \end{scope}
    \node[titlelabel] at (layer3.north west) [yshift=0.25cm] {\textbf{Layer 3: PANOPTES Multimodal Causal Correlation \& Diagnostic Reasoning Engine}};

    % =========================================================================
    % LAYER 4: CLOSED-LOOP REMEDIATION & VERIFICATION
    % =========================================================================
    \node[actnode] (act_operator) at (0.0, -0.65) {
        \textbf{Autonomous Diagnostic Verification \& Kubernetes Operator}\\
        Standardized JSONL Evidence Audit Trail $\;|\;$ Automated Container Rescheduling $\;|\;$ SLO Verification
    };

    \begin{scope}[on background layer]
        \node[layerbox, fit=(act_operator)] (layer4) {};
    \end{scope}
    \node[titlelabel] at (layer4.north west) [yshift=0.25cm] {\textbf{Layer 4: Closed-Loop Feedback \& Autonomous Diagnostic Remediation}};

    % =========================================================================
    % CONNECTIONS & DATA FLOWS (ZERO CROSSING / ZERO OVERLAP)
    % =========================================================================
    % Layer 1 Intra: SUT to Probes (Straight vertical lines)
    \draw[arrow] (sut_frontend) -- (probe_traces);
    \draw[arrow] (sut_shipping) -- (probe_metrics);
    \draw[arrow] (sut_payment)  -- (probe_logs);

    % Layer 1 to Layer 2: Probes to Storage (Straight parallel vertical lines)
    \draw[arrow] (probe_traces.south)  -- node[right, font=\sffamily\tiny] {OTLP Spans} (store_tempo.north);
    \draw[arrow] (probe_metrics.south) -- node[right, font=\sffamily\tiny] {Prom Scrapes} (store_prom.north);
    \draw[arrow] (probe_logs.south)    -- node[right, font=\sffamily\tiny] {Log Streams} (store_loki.north);

    % Layer 2 Intra: OpenMetrics Exemplar Bridge (Horizontal adjacent line between Tempo & Prometheus)
    \draw[biarrow, draw=teal!85!black] (store_tempo.east) -- node[above, font=\sffamily\tiny\bfseries] {OpenMetrics Exemplars} (store_prom.west);

    % Layer 2 to Layer 3: Storage to Correlation Algorithms (Straight parallel vertical lines)
    \draw[arrow, purple!85!black] (store_tempo.south) -- node[left, font=\sffamily\tiny] {Trace Trees} (algo_selftime.north);
    \draw[arrow, purple!85!black] (store_prom.south)  -- node[right, font=\sffamily\tiny] {PromQL $\Delta M > 3\sigma$} ([xshift=-2.8cm]algo_fusion.north);
    \draw[arrow, purple!85!black] (store_loki.south)  -- node[right, font=\sffamily\tiny] {LogQL Crash Sigs} ([xshift=2.8cm]algo_fusion.north);

    % Layer 3 Intra: Algorithm 1 to Algorithm 2 (Straight horizontal connection)
    \draw[arrow, purple!85!black] (algo_selftime.east) -- node[above, font=\sffamily\tiny\bfseries] {Disambiguated $T_{\text{self}}$} (algo_fusion.west);

    % Layer 3 to Layer 4: Authoritative Verdict (Straight vertical down)
    \draw[arrow, orange!90!black] ([xshift=-2.8cm]algo_fusion.south) -- node[right, font=\sffamily\tiny\bfseries] {Authoritative Verdict $\mathcal{V}^*$} (act_operator.north);

    % Layer 4 Feedback Loop to Layer 1 (Clean external bypass route on the left margin)
    \draw[arrow, blue!75!black, dashed] (act_operator.west) -- ++(-2.4, 0) |- node[near end, left, font=\sffamily\tiny, align=right] {Dynamic Telemetry Throttling\\\& Automated Healing Loop} ([yshift=-0.1cm]layer1.west);

\end{tikzpicture}
\caption{The PANOPTES Reference Architecture: A 4-layer framework integrating non-invasive kernel/user probing, federated multimodal storage with OpenMetrics Exemplars, span self-time decomposition with cross-modal evidence fusion, and closed-loop diagnostic remediation.}
\Description{Diagram illustrating the four structural layers of the PANOPTES reference architecture: Ingestion and Probing, Federated Multimodal Storage, Causal Correlation Engine, and Closed-Loop Remediation.}
\label{fig:panoptes_architecture}
\end{figure*}

\begin{table*}[!htbp]
\centering
\small
\caption{Architectural Taxonomy and Structural Layers of the PANOPTES Reference Architecture}
\label{tab:panoptes_layers}
\begin{tabularx}{\textwidth}{@{}l p{3.6cm} p{4.2cm} X@{}}
\toprule
\textbf{Layer} & \textbf{Core Responsibilities} & \textbf{Underlying Technologies} & \textbf{Semantic Invariants \& Bindings} \\ \midrule
\textbf{Layer 1: Probing \& Ingestion} & Non-invasive telemetry acquisition at user and kernel space & eBPF (Grafana Beyla), OpenTelemetry Collector DaemonSet, Promtail & Context propagation via W3C TraceContext; kernel probes capture L4/L7 socket events without application bytecode recompilation. \\ \addlinespace[2pt]
\textbf{Layer 2: Multimodal Storage} & High-throughput, federated persistence of metrics, logs, traces & Prometheus (TSDB), Grafana Loki (Chunk store), Grafana Tempo & OpenMetrics Exemplar binding linking high-cardinality span IDs directly to metric time series buckets; label namespace consistency. \\ \addlinespace[2pt]
\textbf{Layer 3: Causal Correlation} & Disambiguating cascading delays from true root-cause anomalies & PANOPTES Correlation Engine (PromQL, LogQL, Tempo gRPC API) & Exclusive Span Self-Time decomposition ($T_{\text{self}} = T_{\text{span}} - \sum T_{\text{child}}$); metric saturation delta thresholding ($\Delta M > 3\sigma$); log exception signature clustering. \\ \addlinespace[2pt]
\textbf{Layer 4: Feedback \& Remediation} & Closed-loop actuation, SLO reporting, and automated healing & Kubernetes Custom Controllers, Chaos Mesh / Litmus API integration & Closed-loop verification; dynamic probe frequency dampening during cluster distress; automated alerting routing to incident management. \\ \bottomrule
\end{tabularx}
\end{table*}


\subsection{Layer 1: Telemetry Ingestion \& Non-Invasive Probes}
Layer 1 is responsible for capturing raw runtime signals from both application user space and the underlying Linux operating system kernel without imposing unmanageable performance penalties:
\begin{itemize}[leftmargin=1.2em, itemsep=2pt, topsep=3pt]
    \item \textbf{Non-Invasive eBPF Instrumentation (Grafana Beyla):} To observe network sockets, HTTP/2 payloads, and gRPC remote procedure calls without bytecode modification or library re-linking, PANOPTES deploys kernel-resident extended Berkeley Packet Filter (eBPF) probes. Attached to tracepoints such as \texttt{kprobe/tcp\_v4\_connect} and \texttt{uprobe/ssl\_write}, eBPF programs calculate round-trip socket latencies and export standardized OpenTelemetry spans with zero application downtime.
    \item \textbf{OpenTelemetry Collector DaemonSet:} Standardized OTLP (OpenTelemetry Protocol) pipelines run as a lightweight per-node DaemonSet, receiving traces and metrics via gRPC, executing local batching, memory-limiter throttling, and attribute normalization (e.g., standardizing \texttt{k8s.pod.name} and \texttt{k8s.namespace.name} labels) prior to network egress.
    \item \textbf{Linux cgroup v2 Resource Scrapers:} Node-level resource usage is captured directly from the Linux unified \texttt{cgroup} filesystem by the Kubernetes Kubelet, ensuring authoritative metric accounting decoupled from container runtime telemetry agents.
\end{itemize}

\subsection{Layer 2: Multimodal Federated Telemetry Backends}
Rather than enforcing an unwieldy, monolithic database schema that struggles to balance high-throughput time-series metrics with high-volume span graphs and unstructured log lines, Layer 2 employs a federated persistence model governed by three specialized, best-of-breed open-source backends:
\begin{itemize}[leftmargin=1.2em, itemsep=2pt, topsep=3pt]
    \item \textbf{Prometheus (Time-Series Metrics TSDB):} Optimized for high-frequency numeric scraping ($15\,\text{s}$ scrape interval), sub-millisecond gauge lookups, and counter delta evaluations via PromQL.
    \item \textbf{Grafana Loki (Log Stream Store):} Inverts conventional full-text indexing by organizing log streams strictly by Kubernetes metadata labels, allowing streaming regex searches and exception frequency aggregations via LogQL without storage bloat.
    \item \textbf{Grafana Tempo (Distributed Trace Object Store):} Implements massively scalable, index-free object storage for distributed traces, indexing solely by 128-bit \texttt{trace\_id}.
\end{itemize}

\subsubsection{The Semantic Binding: OpenMetrics Exemplar Bridge}
The fundamental innovation enabling cross-modal cohesion between Layer 1 and Layer 2 is the \textbf{OpenMetrics Exemplar Bridge}. When the OpenTelemetry Collector processes an incoming request span, it attaches an Exemplar—a discrete metadata tuple containing $(\texttt{trace\_id}, \texttt{span\_id}, \text{timestamp})$—directly to the corresponding Prometheus histogram bucket. When a latency spike is detected in a PromQL metric query, the Exemplar serves as an immediate pointer, allowing Layer 3 to jump directly to the representative distributed trace without requiring exhaustive cross-database full scans.

\subsection{Layer 3: Causal Correlation \& Diagnostic Reasoning Engine}
Layer 3 represents the algorithmic core of PANOPTES. It eliminates the manual dashboard correlation previously required of human SREs by fusing telemetry streams into an authoritative diagnostic verdict:
\begin{itemize}[leftmargin=1.2em, itemsep=2pt, topsep=3pt]
    \item \textbf{Span Self-Time Decomposition:} As formally detailed in Section~\ref{sec:algorithms}, this module subtracts the cumulative duration of all downstream child spans from parent spans, neutralizing the Cascading Span Bias and accurately isolating the true RPC bottleneck.
    \item \textbf{PromQL Resource Saturation Detection:} Continuously samples cgroup CPU throttling and memory working sets, triggering when usage deviates beyond a dynamic threshold ($\Delta M > 3\sigma_{\text{baseline}}$).
    \item \textbf{LogQL Crash \& Exception Pattern Matching:} Filters log streams for critical container lifecycle signatures (e.g., \texttt{SIGTERM}, \texttt{SIGKILL}, panic backtraces, connection-refused exceptions) within the temporal window of interest.
\end{itemize}

\subsection{Layer 4: Closed-Loop Feedback \& Autonomous Remediation}
The final layer operationalizes the diagnostic outputs of Layer 3 to close the control loop:
\begin{itemize}[leftmargin=1.2em, itemsep=2pt, topsep=3pt]
    \item \textbf{Autonomous Diagnostic Verification:} Exposes REST and gRPC endpoints that emit standardized JSON diagnostic audit trails, logging the identified root-cause microservice, fault modality, and supporting evidence IDs (e.g., specific PromQL vectors, Loki log excerpts, or Tempo span identifiers).
    \item \textbf{Kubernetes Custom Controller Integration:} Emits structured Kubernetes events or Custom Resource updates, allowing automated auto-scalers, restart controllers, or circuit breakers to execute remediation without human latency.
    \item \textbf{Dynamic Telemetry Throttling:} When cluster resource pressure exceeds established safety budgets ($>80\%$ node CPU), Layer 4 dispatches feedback signals to Layer 1, adaptively reducing trace sampling rates and metric scraping frequencies to protect workload availability.
\end{itemize}
